
\chapter*{\RL{ملخص}}
\addcontentsline{toc}{chapter}{Arabic Abstract}

\begin{RLtext}

مع التطور المتسارع لتقنيات الذكاء الاصطناعي، أثبتت النماذج اللغوية قدرة وكفاءة عالية جدا في استيعاب المعلومات الضخمة وإعادة صياغتها للإجابة على التساؤلات المعقدة. في المقابل، لا تزال الاستشارات القانونية التقليدية وتحليل الملفات عملية مكلفة للغاية. في حين أن التوظيف السليم للذكاء الاصطناعي يمكن أن يقلص هذه التكلفة بشكل جذري لتصبح تكلفة شبه رمزية، مما يسهل ولوج الأفراد للخدمات القانونية.

هدفنا من خلال مشروع نهاية السنة (PFA) هو تطوير منصة رقمية متكاملة تتيح للمستخدمين الحصول على استشارات قانونية دقيقة، تحليل ملفاتهم المعقدة، وصياغة وتوليد الوثائق القانونية بشكل آلي. لتحقيق هذه الأهداف، قمنا بدايةً بدراسة احتياجات السوق وتحديد الوظائف الأساسية اللازمة لحل الإشكاليات المطروحة، ثم قمنا بنمذجة النظام وهيكلته الشاملة باستخدام مخططات UML.

واجهتنا عقبة جوهرية تتمثل في أن النماذج اللغوية العالمية المعروفة (مثل Gemini و GPT و Claude) لم يتم تدريبها بشكل كافٍ على النصوص التشريعية المغربية، مما يجعلها عرضة لظاهرة "الهلوسة" واختلاق معلومات خاطئة. لحل هذه المشكلة، قمنا بهندسة وبناء هيكلية متقدمة تعتمد على تقنية الاسترجاع المعزز بالتوليد (RAG). هذه الهيكلية مكنتنا من بناء قاعدة بيانات متجهة باستخدام تقنية ChromaDB، معتمدين في تغذيتها على تقنية التعرف الضوئي على الحروف (OCR) عبر نموذج GPT-4o، مع إدخال تحسينات مصممة خصيصاً لتلائم الوثائق المغربية.

أما على مستوى البنية التحتية، فقد اعتمدنا على إطار العمل FastAPI لتطوير الواجهة الخلفية، ومكتبة React لبناء واجهات المستخدم، في حين استخدمنا نظام PostgreSQL لإدارة قواعد البيانات.

\vspace{0.5cm}

\textbf{الكلمات المفتاحية:} الاستشارات القانونية، الذكاء الاصطناعي، الهلوسة، RAG، قاعدة البيانات المتجهة، ChromaDB، GPT-4o، تقنية OCR، مخططات UML، FastAPI، React، PostgreSQL.

\end{RLtext}